% !TEX root = ../main.tex
\chapter*{Abstract}
\addcontentsline{toc}{chapter}{Abstract}

Constant false-alarm rate (CFAR) detection remains a standard approach to radar target detection,
while recent work increasingly reports improvements from learned detectors. Such comparisons can be
difficult to interpret when methods observe different inputs, use different spatial context, or are
evaluated at unmatched false-alarm rates. This work compares a lightweight
\TinyParams{}-parameter convolutional detector (TinyCNN) with CA-, GO-, SO- and OS-CFAR on
simulated FMCW range--Doppler maps using comparable $\TinyRFR \times \TinyRFD$ local context and a
shared downstream evaluation pipeline. Performance is compared at a matched measured false-alarm
rate across isolated, multi-target, masking, clutter-edge, mixed and homogeneous-clutter conditions,
with an additional unresolved-target condition used as a resolution and protocol stress test.

On a fresh held-out synthetic test set, the network establishes statistically significant gains
across the five resolvable challenge conditions, while homogeneous clutter is a practical tie and
the unresolved stress test shows no significant difference after multiple-comparison correction. The
network reaches $\Pd = 0.5$ at a signal-to-noise ratio \SI{\SynSnrGainDL}{\deci\bel} below CA-CFAR.

Attribution experiments show that the learned advantage is not explained by simple mainlobe-energy
integration and that access to guard-region information modestly affects the network response. They
also show that localisation behaviour is strongly influenced by supervision design: replacing the
original flat $3\times3$ label with more spatially concentrated supervision substantially improves
localisation with negligible impact on tolerant detection performance.

Finally, zero-shot evaluation on the CARRADA automotive radar dataset, without retraining or
real-data tuning, finds no statistically significant difference between TinyCNN and CA-CFAR at exact
annotated support, while the network produces approximately \RoneBurdenRatio{} times more unmatched
detections per frame. These results show that a lightweight learned detector can provide meaningful
advantages over CFAR under controlled synthetic conditions, but that those gains do not transfer
cleanly to real radar measurements, highlighting the importance of evaluation protocol, supervision
design and domain generalisation.

\vspace{0.4em}
\noindent\textbf{Keywords:} FMCW radar, CFAR detection, convolutional neural networks, domain
shift.
